\Chapter{Cel i zakres projektu}
\label{ch:intro}

Zjazd do alei serwisowej zmienia przebieg wyścigu: samochód traci czas
na przejazd i obsługę, a po wymianie opon może uzyskać inne tempo.
Celem projektu jest oszacowanie na końcu okrążenia kierowcy
prawdopodobieństwa, że na następnym okrążeniu zjedzie on po nowe opony.
System ma łączyć odtwarzalne przygotowanie danych, porównanie metod
uczenia maszynowego i czytelną prezentację prognoz w panelu wyścigu.
Kryterium powodzenia stanowi rzetelne porównanie metod na wspólnych
warunkach, bez założonej z góry przewagi sieci neuronowych.

Jednostką analizy jest próbka $K=(r,d,t)$: wyścig $r$, numer kierowcy $d$
i numer właśnie zakończonego okrążenia $t$. Odpowiedź
$y_{r,d,t}\in\{0,1\}$ ma oznaczać wymianę opon na okrążeniu $t+1$.
Wejście modelu musi odpowiadać wiedzy dostępnej w chwili
$T=\mathrm{Time}(r,d,t)$, czyli na końcu okrążenia $t$. W historycznym
zapisie dalsze okrążenia mogą służyć do ustalenia odpowiedzi, lecz nie
mogą wejść do prognozy. Rozdział~\ref{ch:data} przedstawia konsekwencje
tej zasady dla etykiet, historii kierowców i danych całej stawki.

Przewidywanie zachowania zespołu nie rozstrzyga, czy jego decyzja była
optymalna. Projekt nie wyznacza strategii minimalizującej czas wyścigu
ani nie szacuje skutku alternatywnego momentu zjazdu. Jego zastosowaniem
może być analiza historyczna: zestawienie oczekiwanego zdarzenia
z rzeczywistym przebiegiem wyścigu i warunkami, w których pojawił się
alarm. Ocena takiej użyteczności wymaga późniejszych eksperymentów.

\section{Zakres bieżącego opracowania}

Raport opisuje przyjętą rewizję kodu \code{5ba376f} z 9 listopada
2026 r.~\cite{project}. Obecny rezultat dotyczy infrastruktury danych:
odczytu i kontroli źródła, utrwalania wspólnej wersji wejścia oraz
sprawdzania interfejsów między etapami. W rozdziałach o realizacji
i eksperymentach oddzielono działające mechanizmy od krótkiego planu
ich dalszego wykorzystania. Stopień realizacji całego celu zestawiono
w tabeli~\ref{tab:status}.

Dane obejmują sesje Race sezonów 2022--2025. Ograniczenie do wyścigów
suchych i bez czerwonej flagi pozwala skupić się na węższym, możliwym
do objaśnienia zadaniu. Nie uzasadnia ono przenoszenia wniosków na
sprinty, deszcz lub inne sezony. Kolejne części dokumentu prowadzą od
znaczenia obserwacji przez przepływ danych do oceny tego, co dotychczasowe
kontrole rzeczywiście potwierdzają.
